\section{Data}
\label{sec:data}

\paragraph{Dataset.}
All experiments use \textbf{CIFAR-10}~\cite{krizhevsky2009cifar}, a standard benchmark for image classification consisting of 60{,}000 colour images ($32\times32$ pixels) evenly distributed across 10 classes (airplane, automobile, bird, cat, deer, dog, frog, horse, ship, truck).
The standard split provides 50{,}000 training images and 10{,}000 test images.

\paragraph{Train/validation split.}
We hold out 10\% of the training set (5{,}000 images) as a validation set for checkpoint selection, leaving 45{,}000 images for training.
The split is determined by a fixed random seed and is identical across all experiments so that validation accuracy is directly comparable.

\paragraph{Preprocessing.}
Raw pixel values in $\{0,\ldots,255\}$ are normalised to $[0,1]$ and then standardised channel-wise using the CIFAR-10 training-set statistics ($\mu = (0.4914, 0.4822, 0.4465)$, $\sigma = (0.2470, 0.2435, 0.2616)$).

\paragraph{Augmentation.}
The two training recipes differ in the augmentation applied during training (validation and test images are always clean):

\begin{itemize}[leftmargin=*,topsep=2pt,itemsep=1pt]
    \item \textbf{Weak recipe:} random $32\times32$ crop from a 4-pixel reflect-padded image, random horizontal flip.
    \item \textbf{Strong recipe:} same crop and flip, followed by RandAugment ($n{=}2$, $m{=}9$), then Mixup ($\alpha{=}0.8$) or CutMix ($\alpha{=}1.0$) applied at the batch level (each with 50\% probability per batch).
\end{itemize}

The key difference between recipes is that Mixup and CutMix blend pairs of images at the pixel level, fundamentally changing the distribution seen by the patch embedding.
This is the mechanism behind the whitening reversal documented in Section~\ref{sec:experiments}: the whitened basis is computed from clean images and therefore mismatches the blended inputs.

\paragraph{Whitening basis.}
For experiments with frozen patch-whitening initialisation, the patch covariance matrix is computed from the first 5{,}000 training images (clean, before any augmentation), using all overlapping $4\times4$ patches.
This is identical to the procedure in the original CNN paper~\cite{garg2024cifar10}.
